\chapter{The forecast case}\label{ch:case}

A model like WRF is a large set of options; a forecast uses a small subset of them. The port, and most of this book,
is defined by one case: a coupled weather and fire simulation of the Eaton Fire, which started in the foothills above Altadena,
California, on the evening of 7~January~2025 local time. The case's files are in \code{cases/eaton_20250108/}. The ignition used
in the case is at 34.18604~N, 118.09325~W, at 02:18~UTC on 8~January~2025.

\section{Domains and time steps}

Two domains with one-way nesting (\code{feedback = 0}). The ratio is 9 in both space and time.

\begin{center}
\begin{tabular}{lrrrrr}
\toprule
Domain & Grid & $\Delta x$ & $\Delta t$ & Steps in 17 h & Fire grid \\
\midrule
d01 & $450\times450\times60$ & 900 m & 3 s & 20{,}400 & -- \\
d02 & $811\times811\times60$ & 100 m & 1/3 s & 183{,}600 & $3244\times3244$ (25 m) \\
\bottomrule
\end{tabular}
\end{center}

Each large time step has three Runge--Kutta stages, with 1, 2 and 4 acoustic sub-steps
(\code{time_step_sound = 0} gives four per step): seven acoustic sub-steps in all (\cref{ch:time}).

\section{Options}

\paragraph{Dynamics.}
\begin{itemize}
  \item Non-hydrostatic, with the hybrid vertical coordinate (\code{hybrid_opt = 2}) and moist potential
    temperature (\code{use_theta_m = 1}).
  \item Advection: fifth order in the horizontal, third order in the vertical; positive-definite for moisture and
    TKE.
  \item Turbulence:
    \begin{itemize}
      \item d01: two-dimensional Smagorinsky horizontal diffusion with the YSU boundary layer
        (\code{km_opt = 4});
      \item d02: a three-dimensional TKE large-eddy simulation (\code{km_opt = 2}).
    \end{itemize}
  \item Damping: Rayleigh damping of $w$ at the top (\code{damp_opt = 3}) and vertical-velocity damping
    (\code{w_damping = 1}).
\end{itemize}

\paragraph{Physics.}
\begin{itemize}
  \item WSM6 microphysics (\cref{ch:mp}).
  \item Radiation (\cref{ch:rad}), called every minute:
    \begin{itemize}
      \item RRTMG longwave, in single precision in this build;
      \item Dudhia shortwave.
    \end{itemize}
  \item The revised MM5 surface layer and the Noah land-surface model (\cref{ch:sfc}).
  \item YSU boundary layer on d01 only (\cref{ch:pbl}).
  \item No cumulus, shallow-convection, urban or lake schemes.
\end{itemize}

\paragraph{Fire, on d02.}
\begin{itemize}
  \item WRF-Fire (\code{ifire = 2}), refinement 4 in each direction.
  \item Hybrid WENO5/ENO1 level-set advection (\code{fire_upwinding = 9}), with re-initialization.
  \item Wind interpolated to 1~m with a logarithmic profile.
  \item Heat and moisture fluxes fed back into the atmosphere.
  \item The built-in 53-category fuel table and a uniform fuel moisture of 8\,\%.
  \item One circular ignition of radius 100~m.
\end{itemize}

\section{What the case means for the port}

The options fix the code paths:
\begin{itemize}
  \item every routine the case executes in a time step must run on the GPU;
  \item every routine it does not execute may stay as it is.
\end{itemize}

A startup check in the GPU build refuses a namelist outside the supported envelope. Dates, sizes, the nest position
and the ignition may change; the scheme choices may not. The full list of options, with the reasons for each, is in
\code{plan.md} section~2. \Cref{app:namelist} reproduces the namelist with comments.

\begin{portnote}
For development, two smaller cases are used:
\begin{itemize}
  \item \emph{S-3M}, the WRF \code{em_fire} idealized case with the Eaton physics, runs in minutes on one CPU core.
    It is the case of the CPU-only verification.
  \item \emph{eaton\_small} is the real case with a $181\times181$ nest.
\end{itemize}
\end{portnote}
